\section{Methodology}
\label{sec:methodology}

This section describes the full training pipeline for the \sozkz{} model family:
data collection and cleaning, tokenizer design, model architecture, and training
configuration. All code, configs, and trained models are released publicly to
enable full reproducibility.

% ---------------------------------------------------------------------------
\subsection{Training Data}
\label{sec:data}

We construct a large-scale monolingual Kazakh corpus by aggregating text from
18~web and curated sources, including CulturaX~\citep{conneau2020xlmr},
HPLT~2.0, mC4~\citep{devlin2019bert}, MADLAD-400, mOSCAR, Kazakh Wikipedia,
and the \texttt{kz-transformers/multidomain-kazakh-dataset}.
The raw collection comprises 28.4M~documents.

We apply a 9-stage cleaning pipeline:
\begin{enumerate}
  \item Unicode NFC normalization,
  \item removal of control characters,
  \item whitespace and newline collapsing,
  \item minimum length filtering ($\geq50$~characters),
  \item URL density filtering ($\leq5$ per 1{,}000~characters),
  \item HTML tag filtering ($\leq5$ tags),
  \item Kazakh character ratio filter (script-based),
  \item language identification (retaining only Kazakh),
  \item exact deduplication via MD5 hashing (both within-source and
        cross-source against the existing \texttt{multidomain-kazakh-dataset}
        reference of 12.4M~unique hashes).
\end{enumerate}

After cleaning, 13.7M~documents remain (48.2\% pass rate), yielding
approximately 9.0B~tokens under our BPE~50K tokenizer. This constitutes one of
the largest publicly available monolingual Kazakh corpora. We note the inherent
data-quality versus quantity tradeoff in low-resource settings: aggressive
filtering removes noise but discards potentially useful text, while lenient
filtering risks training on low-quality data. Our 48.2\% pass rate reflects a
moderately conservative stance.

% ---------------------------------------------------------------------------
\subsection{Tokenizer}
\label{sec:tokenizer}

We train a ByteLevel BPE tokenizer~\citep{sennrich2016bpe} with a vocabulary
of 50{,}257~tokens exclusively on Kazakh text from our cleaned corpus.

The motivation for a dedicated tokenizer is twofold. First, multilingual
tokenizers (e.g., those used by Llama, Qwen, or Gemma) allocate vocabulary
capacity across 100+ languages, resulting in poor \emph{fertility}---the average
number of tokens per word---for morphologically rich, Cyrillic-script languages
like Kazakh~\citep{rust2021tokenizer, toraman2023turkic}.
Second, better tokenization directly improves training and inference efficiency:
fewer tokens per document means faster throughput and lower compute cost for a
fixed amount of text.

The tokenizer is published as open-source on HuggingFace.\footnote{\url{https://huggingface.co/stukenov/sozkz-core-gpt2-50k-kk-base-v1}}

% ---------------------------------------------------------------------------
\subsection{Model Architecture}
\label{sec:architecture}

We train four model sizes following the LlamaForCausalLM
architecture~\citep{touvron2023llama}: SwiGLU activations~\citep{shazeer2020glu},
Rotary Position Embeddings (RoPE)~\citep{su2021rope}, RMSNorm
pre-normalization~\citep{zhang2019rmsnorm}, and tied input--output embeddings.
All models are trained from scratch with no pretrained initialization.

\begin{table}[ht]
  \centering
  \caption{Architecture configurations for the \sozkz{} model family.
           All models use SwiGLU, RoPE, RMSNorm, tied embeddings, and
           a vocabulary of 50{,}257 tokens.}
  \label{tab:architecture}
  \begin{tabular}{@{}lrrrrr@{}}
    \toprule
    Model & Params & Layers & Hidden & Heads & Intermediate \\
    \midrule
    \sozkz{}-50M   &  50.3M & 8  & 576   & 8  & 1{,}536 \\
    \sozkz{}-150M  & 151.9M & 16 & 768   & 12 & 2{,}048 \\
    \sozkz{}-300M  & 325M   & 18 & 1{,}024 & 16 & 3{,}584 \\
    \sozkz{}-600M  & 587M   & 22 & 1{,}280 & 20 & 4{,}480 \\
    \bottomrule
  \end{tabular}
\end{table}

The Llama architecture was chosen for its well-understood scaling
properties~\citep{kaplan2020scaling}, training stability, and broad ecosystem
support. The intermediate dimension follows a $3.5\times$ multiplier of the
hidden size across all configurations, consistent with SwiGLU best
practices~\citep{shazeer2020glu}. Context length is 2{,}048~tokens for the
50M, 300M, and 600M~models; the 150M~model uses a 1{,}024-token context.

% ---------------------------------------------------------------------------
\subsection{Training Details}
\label{sec:training}

All models are trained for one epoch over the full ${\sim}9.0$B-token corpus
using the AdamW optimizer~\citep{loshchilov2017adamw} with $\beta_1 = 0.9$,
$\beta_2 = 0.95$, and weight decay of~0.1. We use a cosine learning rate
schedule~\citep{loshchilov2016cosine} with linear warmup. Peak learning rates
are $6 \times 10^{-4}$ for the 50M~model and $3 \times 10^{-4}$ for the 150M
and 600M~models. Gradient clipping is set to~1.0.

Training uses \texttt{bfloat16} mixed precision, which provides better numerical
stability on modern GPUs compared to \texttt{float16}, particularly for the
larger models. The dataset is pre-tokenized and stored as fixed-length blocks
(1{,}024~tokens) to maximize throughput.

For the 600M~model, the data-to-parameter ratio is $9.0\text{B} / 587\text{M}
\approx 15.3{:}1$. This is slightly below the Chinchilla-optimal ratio of
${\sim}20{:}1$~\citep{hoffmann2022chinchilla}, meaning the model is mildly
undertrained relative to the compute-optimal frontier. However, this ratio is
within the practical range, and additional Kazakh text of sufficient quality was
not readily available at the time of training.

\paragraph{Hardware.}
The 50M and 150M~models are trained on $2\times$ NVIDIA RTX~4090 GPUs via
vast.ai using DDP (Distributed Data Parallel). The 600M~model is trained on
$8\times$ NVIDIA H100 SXM 80\,GB GPUs, also on vast.ai, achieving
approximately 577K~tokens/s throughput.

\paragraph{Reproducibility.}
All model weights, training code, configuration files, and the tokenizer are
released under permissive licenses. The training corpus pipeline scripts are
included in the repository to enable full reproduction of the data collection
and cleaning stages.
